\documentclass[11pt]{article}
\usepackage{amsmath,amssymb,mathtools}
\usepackage{fontspec,unicode-math}
\setmainfont{Latin Modern Roman}
\setsansfont{Latin Modern Sans}
\setmathfont{Latin Modern Math}
\usepackage{xurl,hyperref}
\newcommand{\sref}[2]{\hyperlink{source:#1:#2}{[S#2]}}
\newcommand{\cataloguescope}[1]{\paragraph{Mathematical question.}#1}
\begin{document}
\section{Simple Linear-Time Polygon Triangulation}\label{statement}

\subsection{Definitions and mathematical statement}

\paragraph{Shared notation.}$\mathbb N=\{1,2,\ldots\}$, and $\mathbb Z,\mathbb Q,\mathbb R,\mathbb C$ denote the integers, rationals, reals and complex numbers. Any other conventions are specified locally.


A simple polygon has a non-self-intersecting polygonal boundary with $n$ vertices. A triangulation partitions its interior into triangles by noncrossing diagonals, using its vertices. Linear time means $O(n)$ running time. Simplicity of the algorithm is the qualitative criterion used in the source.
\paragraph{Statement recorded in the supplied source.}
Find a deterministic linear-time triangulation algorithm for arbitrary simple polygons that is significantly simpler than Chazelle\textquotesingle s. The related question recorded with it asks for a simple linear-time algorithm computing a shortest path within a simple polygon without first using a more complicated triangulation algorithm. \sref{statement}{2}
\subsection{Short English statement}
Find a deterministic linear-time triangulation algorithm for arbitrary simple polygons that is significantly simpler than Chazelle\textquotesingle s. The related question recorded with it asks for a simple linear-time algorithm computing a shortest path within a simple polygon without first using a more complicated triangulation algorithm.
\subsection{Sources}
\begin{description}
\item[\hypertarget{source:statement:1}{S1}] ULAM.AI and the UnsolvedMath Contributors, UnsolvedMath: A Curated Collection of Open Mathematics Problems, catalogue entry. CC BY 4.0. \url{https://huggingface.co/datasets/ulamai/UnsolvedMath}.
\item[\hypertarget{source:statement:2}{S2}] Erik D. Demaine, Joseph S. B. Mitchell and Joseph O\textquotesingle Rourke (editors), \emph{The Open Problems Project}, maintained original source, Problem 10, Statement and complete Partial/Related Results. \url{https://github.com/edemaine/topp/blob/main/Problems/P.000010}.
\end{description}
\label{end:statement}
\end{document}
